% Tier C structural overview. First float of the Method section, placed before
% the section's first equation. Carries no magnitudes: it states what is fixed,
% what varies, and what the paper reports.
\begin{figure*}[t]
  \centering
  \resizebox{0.99\textwidth}{!}{%
  \begin{tikzpicture}[x=1cm,y=1cm]
    \node[acbox,minimum width=3.5cm,minimum height=1.15cm,align=center] (a) at (1.9,1.1)
      {\textbf{Frozen \evalsem}\\Same bound code\\Same evaluation precision};
    \node[acbox,minimum width=3.5cm,minimum height=1.15cm,align=center] (b) at (6.7,1.1)
      {\textbf{Independent factors}\\\splitseed\\\trajseed};
    \node[acbox,minimum width=3.7cm,minimum height=1.15cm,align=center] (c) at (11.7,1.1)
      {\textbf{Comparable outputs}\\Mean, standard deviation,\\range, paired contrasts};
    \draw[acarrow] (a) -- (b);
    \draw[acarrow] (b) -- (c);
    % text width is a cap, not a floor: at this 4.1cm center spacing the two
    % callouts can grow only downward (more lines), never sideways into each
    % other, regardless of how long the caption text is.
    \node[acsoft,text width=3.6cm,align=center] at (5.0,2.55)
      {The selection method is the treatment, not a source of hidden protocol
       drift};
    \node[acsoft,text width=3.6cm,align=center] at (9.1,2.55)
      {Every reported statistic traces to preserved primary runs};
  \end{tikzpicture}}
  \caption{\textbf{Study design.} The evaluation is fixed before any comparison
    begins, the two sources of randomness are varied independently, and the
    reported statistics are defined in advance. Fixing the left panel is what
    makes the right panel a measurement of the method rather than of the
    protocol.}
  \label{fig:overview}
\end{figure*}
